\chapter{АНАЛИЗ ПРЕДМЕТНОЙ ОБЛАСТИ}
\label{ch:problem_statement}

\section{Особенности предметной области и постановка задачи}

Системы охранного видеонаблюдения охраняемых зон функционируют в условиях непрерывного поступления видеоданных, неоднородности сцен наблюдения и высоких требований к времени реакции на инциденты. В отличие от задач автономного анализа отдельных изображений, в рассматриваемой предметной области необходимо не только обнаружить объект в одном кадре, но и обеспечить его устойчивое сопровождение во времени, интерпретировать поведение в контексте контролируемой зоны и сформировать событие, пригодное для дальнейшего использования внешними подсистемами безопасности \cite{redmon2016yolo,wojke2017deepsort,transactional_outbox}.

Для платформы видеоаналитики реального времени характерны следующие особенности:
\begin{itemize}
    \item многоканальность наблюдения, при которой одновременно обрабатывается несколько видеопотоков от территориально распределенных камер;
    \item изменчивость условий съемки, обусловленная временем суток, погодными факторами, различиями освещенности и появлением помех в кадре;
    \item наличие объектов разных масштабов и скоростей движения, включая людей, транспортные средства и малогабаритные динамичные объекты, в том числе БПЛА;
    \item необходимость различать отдельные факты появления объектов и значимые для оператора события, например пересечение рубежа, проникновение в запретную зону, длительное нахождение в области контроля или появление объекта заданного класса;
    \item наличие ограничений по вычислительным ресурсам, пропускной способности каналов передачи данных и допустимой задержке выдачи результата.
\end{itemize}

С учетом указанных особенностей задачу разработки платформы можно представить следующим образом. Пусть на вход системы поступает множество видеопотоков
\begin{equation}
\mathcal{V} = \{V_i\}_{i=1}^{N},
\end{equation}
\where
\whereitem{\(V_i\)}{видеопоток камеры \(i\)}
\whereitem{\(x_t^i\)}{последовательность кадров потока \(V_i\)}
\whereitem{\(m_i\)}{метаданные камеры}
\whereitemlast{\(c_i\)}{конфигурационные параметры: зоны, рубежи, классы и правила генерации событий}

Для каждого кадра $x_t^i$ необходимо определить множество детекций
\begin{equation}
\mathcal{D}_t^i = \{d_k\},
\end{equation}
\where
\whereitemlast{\(d_k\)}{элемент множества детекций, содержащий положение объекта в кадре, его класс и оценку достоверности}
На основе последовательности детекций необходимо формировать множество траекторий
\begin{equation}
\mathcal{T}_t^i = \{\tau_j\},
\end{equation}
сохраняющих идентичность объекта на протяжении времени наблюдения. Далее по множеству траекторий, правилам предметной области и описаниям зон контроля требуется сформировать множество событий
\begin{equation}
\mathcal{E}_t^i = \{e_r\},
\end{equation}
каждое из которых характеризуется типом события, временем возникновения, привязкой к камере, связанными объектами и уровнем приоритета.

Таким образом, постановка задачи включает четыре взаимосвязанных этапа:
\begin{enumerate}
    \item прием и предварительную обработку видеопотоков;
    \item обнаружение и классификацию объектов в отдельных кадрах;
    \item сопровождение объектов между кадрами с сохранением идентификаторов;
    \item событийный анализ траекторий и формирование сведений о событиях для внешних потребителей.
\end{enumerate}

Результатом работы платформы должны являться не только детекции и траектории, но и структурированные сообщения о событиях, пригодные для передачи в операторские интерфейсы, журналы регистрации, системы оповещения и другие подсистемы безопасности. При этом система должна работать в режиме, близком к реальному времени, обеспечивая приемлемый баланс между точностью аналитики, устойчивостью сопровождения, скоростью обработки и надежностью доставки событий.

\subsection{Пользовательские роли и сценарии использования}
\label{subsec:user_journey}

Помимо технической постановки задачи, для разработки платформы важно зафиксировать целевые роли пользователей и типовые сценарии их работы с системой. Это позволяет уже на этапе анализа предметной области согласовать функциональные требования с реальной деятельностью операторов и администраторов.

В платформе выделены три пользовательские роли:
\begin{itemize}
    \item \textit{Оператор охраны} (\texttt{operator}) - сотрудник, осуществляющий непрерывный мониторинг охраняемых зон, реагирующий на поступающие тревоги и регистрирующий результаты обработки инцидентов;
    \item \textit{Администратор системы} (\texttt{admin}) - сотрудник, отвечающий за конфигурацию камер, описание зон, настройку правил событийного анализа и контроль учетных записей;
    \item \textit{Аналитик безопасности} (\texttt{viewer}) - сотрудник, работающий с журналом событий и аналитическими отчетами без права изменения конфигурации.
\end{itemize}

Карта пути пользователя (Customer Journey Map, CJM) для роли оператора охраны включает следующие основные этапы:
\begin{enumerate}
    \item \textit{Начало смены.} Оператор выполняет вход в систему, проверяет состояние камер и активных правил, убеждается в исправности интеграционного контура (доставка тестовых событий).
    \item \textit{Непрерывный мониторинг.} На рабочем месте отображаются текущие состояния камер, сводный поток событий и индикаторы операционных метрик. Оператор контролирует ход обработки и при необходимости фильтрует поток по типу события или зоне.
    \item \textit{Реакция на событие.} При поступлении тревоги оператор открывает карточку события, оценивает контекст (тип правила, камера, зона, объект, оценка достоверности), при необходимости просматривает связанный фрагмент видео и фиксирует результат обработки.
    \item \textit{Эскалация и взаимодействие.} В случае подтверждения инцидента событие передается во внешние подсистемы оповещения и управления безопасностью; ход эскалации сохраняется в журнале для последующего разбора.
    \item \textit{Завершение смены.} Оператор формирует сводку обработанных событий, передает изменения сменщику и фиксирует наблюдения по работе платформы (сбои, ложные срабатывания, замечания к настройкам правил).
\end{enumerate}

Аналогичные карты составлены для администратора системы (управление камерами, зонами, правилами и учетными записями) и аналитика безопасности (работа с историей событий, формирование отчетов). Совокупность ролей и сценариев непосредственно используется далее при формировании функциональных и нефункциональных требований к платформе.

\section{Обзор существующих решений и литературный обзор}
\label{sec:related_work}

Согласно методическим рекомендациям по содержанию первой главы выпускной квалификационной работы, на этапе анализа предметной области необходимо рассмотреть существующие подходы и программные платформы, выявить их сильные и слабые стороны и тем самым обосновать актуальность собственной разработки.

\subsection{Краткий литературный обзор}

В научной литературе задачи обнаружения и сопровождения объектов в видеопотоках рассматриваются в значительном числе работ. Алгоритмы одноэтапной детекции класса YOLO впервые предложены в работе~\cite{redmon2016yolo} и в дальнейшем развиты в семействе моделей различной вычислительной сложности~\cite{redmon2018yolov3,ultralytics_docs}, что обеспечило их широкое практическое применение для потоковой обработки видео. Для предметной области БПЛА актуальны специализированные наборы данных и эталонные наборы испытаний: например, UAVDB фиксирует проблему малых размеров объекта, вариативности масштаба и необходимости слабоконтролируемой разметки для высокоразрешающих сцен наблюдения~\cite{chen2026uavdb}. Двухэтапные методы и трансформерные архитектуры обеспечивают более высокое качество распознавания на сложных сценах, однако требуют существенно больших вычислительных ресурсов и уступают одноэтапным подходам по показателю задержки.

В области сопровождения объектов в течение последних лет получили распространение методы сопровождения по результатам обнаружения. Алгоритм SORT~\cite{bewley2016sort} показал, что комбинация фильтра Калмана~\cite{kalman1960} и венгерского алгоритма~\cite{kuhn1955hungarian} позволяет получать высокую скорость сопровождения при умеренной точности на простых сценах. Дальнейшее развитие этого направления привело к алгоритму DeepSORT~\cite{wojke2017deepsort}, дополняющему кинематическое сопоставление признаками внешнего вида и существенно снижающему число переключений идентификаторов. Для оценки качества сопровождения используются методики MOT-Challenge и метрики MOTMetrics~\cite{bernardin2008motmetrics,motchallenge}, а в качестве типовой выборки для предобучения моделей детекции - набор данных COCO~\cite{coco_dataset}.

Построение событийно-ориентированных платформ и надежной асинхронной интеграции хорошо освещено в инженерной литературе. В монографии~\cite{richardson_microservices} и серии онлайн-материалов~\cite{transactional_outbox,rabbitmq_dlq} систематизированы шаблоны микросервисной архитектуры, включая шаблон транзакционной таблицы исходящих сообщений и шаблоны контролируемой повторной доставки. Документация используемых компонентов~\cite{fastapi_docs,sqlalchemy_docs,rabbitmq_docs,rabbitmq_confirms,postgresql_docs,docker_docs,prometheus_docs,grafana_docs} формирует справочную базу для практической реализации платформы.

Нормативную и методическую основу оформления и экспериментальной оценки результатов задают документы системы СИБИД, в том числе ГОСТ~7.32-2017~\cite{gost_732_2017}, ГОСТ~Р~7.0.5-2008~\cite{gost_r_705_2008} и связанные стандарты библиографического описания~\cite{gost_71_2003,gost_780_2000,gost_782_2001}. Для прикладного контура видеонаблюдения также учитываются актуальные международные рекомендации по проектированию, вводу в эксплуатацию и испытанию интеллектуальной видеоаналитики~\cite{iec_62676_4_2025,iec_62676_6_2026}, а для контура безопасности - контрольные семейства доступа и аудита~\cite{nist_sp800_53r5}. Сочетание академической литературы по компьютерному зрению с инженерными подходами к построению распределенных систем и нормативной базой по оформлению научно-технической документации формирует методический контур настоящей работы.

\subsection{Обзор существующих платформ видеоаналитики}

Среди представленных на рынке платформ интеллектуального видеонаблюдения можно условно выделить три основные группы. К первой относятся системы класса VMS (Video Management System) с расширенными модулями видеоаналитики, например \textit{Macroscop}, \textit{TRASSIR} и \textit{Axxon One}; открытые материалы производителей описывают их как платформы видеонаблюдения с развитыми аналитическими модулями~\cite{macroscop_site,trassir_site,axxon_one}. Ко второй группе относятся облачные сервисы видеонаблюдения с аналитикой, например \textit{Ivideon}; они снижают порог входа для пользователя, однако предполагают передачу видеоданных во внешний контур обработки~\cite{ivideon_analytics}. Третью группу образуют международные платформы с акцентом на машинное обучение и расширяемую аналитику, например \textit{Milestone XProtect} и \textit{Vaidio}; такие решения обладают развитой функциональностью, но их открытые материалы не раскрывают внутренний алгоритмический и интеграционный контур в объеме, достаточном для воспроизводимого исследовательского стенда~\cite{milestone_xprotect,vaidio_platform}.

Сравнение характеристик представленных платформ и разрабатываемого решения по существенным для исследования критериям приведено в таблице~\ref{tab:analogues_comparison}. Сравнение выполнено по открытым материалам производителей; если сведения о внутреннем устройстве продукта в них не раскрыты, это фиксируется как ограничение открытой проверяемости, а не как утверждение о фактическом отсутствии соответствующего механизма.

\begin{table}[H]
\caption{Сравнение существующих платформ видеоаналитики и разрабатываемого решения}
\label{tab:analogues_comparison}
\begingroup
\scriptsize
\setlength{\tabcolsep}{2pt}
\begin{tabularx}{\textwidth}{@{}>{\raggedright\arraybackslash}p{0.20\textwidth} >{\raggedright\arraybackslash}p{0.13\textwidth} >{\raggedright\arraybackslash}p{0.13\textwidth} >{\raggedright\arraybackslash}p{0.12\textwidth} >{\raggedright\arraybackslash}p{0.13\textwidth} X@{}}
\toprule
Критерий сравнения & Macroscop & TRASSIR & Ivideon & Vaidio & Платформа \\
\midrule
Контур применения по открытым материалам & локальный/РФ & локальный/РФ & облачный сервис & международный & локальный стенд/РФ \\
Тип развертывания & локальное & локальное & облачное & локальное / облачное & локальное \\
Публично раскрытое алгоритмическое ядро & не раскрыто & не раскрыто & не раскрыто & не раскрыто & открытое YOLO-ядро \\
Возможность замены детектора пользователем & ограничена & ограничена & нет & ограничена & да (по контракту) \\
Многообъектное сопровождение с признаками внешнего вида & частично & да & частично & да & да, DeepSORT-подход \\
Правиловой движок событий (zone\_enter, line\_cross, dwell\_time) & да & да & частично & да & да \\
Надежная асинхронная доставка событий (исходящие сообщения + очередь недоставленных сообщений) & не раскрыта & не раскрыта & облачный контур & не раскрыта & да, на открытом стеке \\
Открытые контракты API/событий & ограничены & ограничены & ограничены & ограничены & да, OpenAPI \\
Прозрачные операционные метрики (Prometheus / Grafana) & не заявлены как открытый контур & не заявлены как открытый контур & нет & не заявлены как открытый контур & да \\
Воспроизводимость стенда (compose/IaC) & не заявлена & не заявлена & не заявлена & не заявлена & да \\
Стоимость владения & лицензия + поддержка & лицензия + поддержка & подписка & лицензия + поддержка & открытый стенд \\
\bottomrule
\end{tabularx}
\endgroup
\end{table}

Анализ показывает, что существующие платформы успешно решают значительную часть прикладных задач охранного видеонаблюдения, однако обладают рядом ограничений, существенных для научно-исследовательской работы и построения предметно-ориентированных решений:
\begin{itemize}
    \item ограниченная раскрытость алгоритмического ядра в открытых материалах затрудняет адаптацию детектора и модуля сопровождения к специфическим условиям конкретного объекта (например, обнаружение БПЛА на дальних дистанциях);
    \item ограниченная открытость контрактов событий и программных интерфейсов (API) усложняет интеграцию с произвольными внешними системами и сравнительные исследования;
    \item ограниченная прозрачность операционных метрик затрудняет независимую оценку эксплуатационных характеристик (сквозной задержки, устойчивости доставки, влияния отказов);
    \item модель развертывания платформ преимущественно ориентирована на эксплуатационное внедрение, а не на воспроизводимый исследовательский стенд.
\end{itemize}

Выявленные ограничения формируют практическую и научную нишу для разрабатываемого решения: открытый алгоритмический и интеграционный контур, фиксированные контракты событий, прозрачные метрики и воспроизводимый стенд развертывания позволяют проводить экспериментальную оценку архитектурных и алгоритмических решений в условиях, сопоставимых с реальными системами охраны, и при этом сохранять контроль за исходным кодом, конфигурацией и наблюдаемостью.

\subsection{Выводы по обзору}

Результаты обзора литературы и существующих платформ позволяют сделать следующие основные выводы:
\begin{enumerate}
    \item для задач охранной видеоаналитики реального времени накоплена развитая алгоритмическая база (YOLO-семейство, SORT/DeepSORT, метрики MOT-Challenge), пригодная для построения базового конвейера обработки;
    \item инженерные шаблоны микросервисной архитектуры (транзакционная таблица исходящих сообщений, контролируемые повторы, очереди недоставленных сообщений) обеспечивают необходимый уровень надежности доставки событий и могут быть перенесены в предметную область охранного видеонаблюдения;
    \item существующие промышленные платформы решают значительную часть задач, однако их закрытость и ограниченная прозрачность контрактов и метрик создают научно-методическую нишу для построения открытого исследовательского решения с фиксированными контрактами и воспроизводимым стендом.
\end{enumerate}

\section{Требования к разрабатываемой платформе}

Формирование требований необходимо для корректного выбора методов и архитектурных решений. Поскольку разрабатываемая система ориентирована на реальную эксплуатацию в составе охранного видеонаблюдения, требования должны охватывать как функциональные возможности, так и ограничения по качеству обслуживания.

\subsection{Функциональные требования}

К основным функциональным требованиям относятся:
\begin{enumerate}
    \item прием одного или нескольких видеопотоков от камер наблюдения и поддержка их непрерывной обработки;
    \item выделение целевых объектов в кадре с определением их пространственного положения и класса;
    \item сопровождение обнаруженных объектов во времени с присвоением устойчивых идентификаторов;
    \item поддержка анализа нескольких классов объектов, включая людей, транспортные средства и БПЛА;
    \item формирование событий на основе пространственно-временных правил, связанных с зонами контроля, виртуальными линиями и логикой поведения объекта;
    \item сохранение результатов обработки в форме метаданных, журналов событий и параметров траекторий;
    \item передача сведений о событиях внешним потребителям через программные интерфейсы или каналы интеграции;
    \item возможность настройки состава правил, перечня контролируемых зон и параметров обработки без изменения базовой логики платформы.
\end{enumerate}

\subsection{Нефункциональные требования}

Для рассматриваемого класса систем определяющее значение имеют следующие нефункциональные требования:
\begin{itemize}
    \item \textit{Работа в реальном времени}. Платформа должна обеспечивать обработку видеопотоков с ограниченной задержкой, достаточной для практического реагирования на инцидент.
    \item \textit{Масштабируемость}. Архитектура должна допускать увеличение числа камер, правил событийного анализа и внешних потребителей без полной переработки системы.
    \item \textit{Отказоустойчивость}. Отдельные сбои обработки, кратковременная недоступность камер или каналов связи не должны приводить к потере работоспособности всей платформы.
    \item \textit{Точность и устойчивость}. Методы должны сохранять приемлемое качество работы при частичных перекрытиях объектов, изменениях освещенности и наличии фоновых помех.
    \item \textit{Расширяемость}. Система должна допускать замену или дообучение детектора, замену модуля сопровождения и уточнение набора событийных правил без кардинальной перестройки конвейера.
    \item \textit{Интегрируемость}. Результаты работы платформы должны представляться в форме, совместимой с внешними системами хранения, мониторинга и оповещения.
    \item \textit{Интерпретируемость}. Для задач безопасности важно не только обнаружить событие, но и объяснить причину его формирования через траекторию объекта, тип правила и контекст наблюдения.
\end{itemize}

Совокупность перечисленных требований определяет необходимость выбора методов, сочетающих высокую скорость обработки с устойчивой работой на реальных видеоданных и возможностью последующей интеграции в событийно-ориентированную программную платформу.

\section{Обоснование выбора методов}

\subsection{Выбор метода обнаружения объектов}

Методы обнаружения объектов, применимые в системах видеонаблюдения, можно условно разделить на несколько групп.

Классические методы, основанные на вычитании фона, межкадровой разности, пороговой обработке и выделении контуров, отличаются сравнительно низкой вычислительной сложностью и могут эффективно использоваться для выявления факта движения при статичной камере. Однако данные методы плохо переносят сложные погодные условия, переменное освещение, тени, дрожание камеры и не обеспечивают устойчивой классификации объектов по типам. Для задач событийного анализа этого недостаточно, поскольку системе необходимо различать классы объектов и корректно обрабатывать сцены с несколькими целями.

Двухэтапные нейросетевые методы обнаружения, например подходы семейства R-CNN, обеспечивают высокую точность локализации и классификации объектов. Их недостатком является относительно высокая вычислительная стоимость, особенно при обработке многоканальных видеопотоков в режиме реального времени. В рассматриваемой задаче чрезмерная задержка обработки снижает практическую ценность решения, даже если точность отдельных детекций остается высокой.

Методы трансформерного типа представляют значительный научный интерес и демонстрируют хорошие результаты в ряде задач компьютерного зрения. Вместе с тем их эффективное применение обычно требует существенных вычислительных ресурсов, значительного объема данных для адаптации и более сложного контура развертывания. Для платформы, ориентированной на эксплуатационные сценарии охранного видеонаблюдения, приоритетным является более предсказуемый баланс между скоростью, точностью и простотой практического внедрения.

С учетом поставленных требований наиболее целесообразным является использование одноэтапного детектора семейства YOLO \cite{redmon2016yolo,redmon2018yolov3,ultralytics_docs}. Выбор данного класса методов обусловлен следующими факторами:
\begin{itemize}
    \item обнаружение и классификация выполняются в едином вычислительном конвейере, что обеспечивает высокую скорость обработки;
    \item методы семейства YOLO хорошо адаптированы к потоковой обработке видео и практическому развертыванию на GPU- и edge-платформах;
    \item использование многомасштабного представления признаков позволяет обнаруживать объекты различных размеров, что существенно для задач наблюдения за протяженными сценами и для сценариев обнаружения БПЛА;
    \item архитектуры данного класса допускают дообучение на предметно-ориентированных наборах данных, что позволяет адаптировать модель к особенностям конкретных охраняемых зон.
\end{itemize}

Следовательно, в качестве базового метода обнаружения в работе целесообразно использовать одноэтапный детектор семейства YOLO, обеспечивающий приемлемое сочетание быстродействия и качества распознавания для задач реального времени.

\subsection{Выбор метода сопровождения объектов}

После обнаружения объектов необходимо обеспечить их сопровождение между кадрами. Изолированная покадровая детекция не позволяет надежно анализировать поведение объекта, определять направление движения и выделять события, зависящие от длительности нахождения в зоне или пересечения рубежей.

Методы, основанные только на оптическом потоке или корреляционных фильтрах, подходят преимущественно для сопровождения отдельных объектов и подвержены накоплению ошибки при длительном наблюдении. Алгоритм SORT обладает высокой скоростью и хорошо подходит для сценариев с качественными детекциями и небольшим числом перекрытий, однако использует в основном кинематические признаки и потому чувствителен к переключениям идентификаторов в сложных сценах.

Для задач охранного видеонаблюдения более обоснованным является использование подхода сопровождения по результатам обнаружения, сочетающего модель движения объекта и процедуру ассоциации новых детекций с уже существующими траекториями. Среди практических методов этого класса целесообразно выделить DeepSORT, который дополняет кинематическое сопоставление признаками внешнего вида объекта \cite{bewley2016sort,wojke2017deepsort}. Это особенно важно при частичных перекрытиях, кратковременных пропаданиях детекций и наличии нескольких близко расположенных целей.

Выбор метода сопровождения типа DeepSORT обусловлен следующими причинами:
\begin{itemize}
    \item использование фильтра Калмана позволяет прогнозировать положение объекта и сглаживать влияние пропусков детекций;
    \item задача ассоциации решается формализованно и допускает устойчивую работу в многообъектных сценах;
    \item применение дескрипторов внешнего вида снижает число переключений идентификаторов по сравнению с чисто кинематическими методами сопровождения;
    \item выход модуля сопровождения естественным образом формирует траектории, необходимые для последующего событийного анализа.
\end{itemize}

Таким образом, для подсистемы сопровождения обоснован выбор многообъектного модуля сопровождения по результатам обнаружения, реализуемого на основе фильтра Калмана, процедуры ассоциации и признаков внешнего вида, что соответствует методологической логике DeepSORT.

\subsection{Выбор метода событийного анализа}

Событийный анализ в системах охранного видеонаблюдения может быть реализован различными способами. Один из подходов предполагает прямое распознавание сложных событий по видеоклипам с помощью специализированных нейросетевых моделей. Такой подход перспективен, однако требует больших объемов размеченных данных именно по событиям, менее прозрачен при интерпретации результата и хуже адаптируется к локальным регламентам конкретного объекта.

Для практико-ориентированной платформы более целесообразным является событийный анализ на основе правил, применяемых к траекториям объектов и конфигурации контролируемых зон. В этом случае событие трактуется как выполнение некоторого пространственно-временного предиката: пересечение виртуальной линии, вход в запретную область, нахождение в зоне дольше порогового интервала, появление объекта заданного класса в определенном секторе наблюдения и т.п. Такой подход лучше согласуется с архитектурным шаблоном надежной асинхронной интеграции через транзакционную таблицу исходящих сообщений \cite{transactional_outbox,richardson_microservices}.

Выбор правилового подхода обусловлен следующими преимуществами:
\begin{itemize}
    \item интерпретируемость результата и возможность явного объяснения причины формирования события;
    \item простота адаптации к конкретным требованиям эксплуатации без полного переобучения модели;
    \item независимость логики событийного анализа от конкретной реализации детектора или модуля сопровождения;
    \item удобство интеграции с подсистемами оповещения, журналирования и операторского контроля.
\end{itemize}

Следовательно, в рамках разрабатываемой платформы рационально использовать композицию из трех уровней: детектор семейства YOLO для покадрового обнаружения, модуль сопровождения на основе DeepSORT и модуль правилового событийного анализа для формирования целевых событий.

\section{Сравнительный анализ альтернатив реализации платформы}

Помимо выбора алгоритмов компьютерного зрения, для практической пригодности системы требуется выбрать технологический стек контура исполнения: способ захвата потоков, модель обмена сообщениями и архитектуру хранения результатов. Ниже приведено сопоставление основных альтернатив.

\subsection{Альтернативы по подсистеме захвата видеопотоков}

\begin{table}[H]
\caption{Сравнение альтернатив захвата видеопотоков}
\label{tab:capture_alternatives}
\begin{tabularx}{\textwidth}{>{\raggedright\arraybackslash}p{0.22\textwidth}X}
\toprule
Альтернатива & Оценка применимости \\
\midrule
\texttt{OpenCV VideoCapture} & Быстрый старт для системы, но ограниченный контроль над буфером, reconnect-политикой и диагностикой транспортного уровня \\
\texttt{GStreamer} & Высокая гибкость, но более сложная конфигурация и порог вхождения на начальном этапе внедрения \\
\texttt{FFmpeg/PyAV} & Компромисс между управляемостью потока, зрелостью экосистемы и интеграцией с Python-контуром \\
\bottomrule
\end{tabularx}
\end{table}

Для текущей работы выбран \texttt{FFmpeg/PyAV}, поскольку это решение дает предсказуемое поведение при сбоях RTSP и не требует усложнения стека на раннем этапе \cite{ffmpeg_docs,pyav_docs}.

\subsection{Альтернативы по интеграционному контуру}

\begin{table}[H]
\caption{Сравнение вариантов интеграции событий}
\label{tab:integration_alternatives}
\begin{tabularx}{\textwidth}{>{\raggedright\arraybackslash}p{0.26\textwidth}X}
\toprule
Вариант & Ограничения/преимущества \\
\midrule
Синхронные HTTP-вызовы из горячего пути & Простая реализация, но высокий риск роста задержки и каскадной деградации при недоступности получателя \\
Kafka-контур на старте & Сильные возможности масштабирования, но избыточная инфраструктурная сложность для диапазона 1--10 камер \\
RabbitMQ + таблица исходящих сообщений & Достаточная надежность доставки, подтверждения публикации, очередь недоставленных сообщений и умеренный порог сопровождения на односерверном стенде \\
\bottomrule
\end{tabularx}
\end{table}

С учетом требований отказоустойчивости и ограниченного масштаба пилотного стенда выбран вариант RabbitMQ + транзакционная таблица исходящих сообщений \cite{rabbitmq_docs,rabbitmq_confirms,transactional_outbox}.

\subsection{Альтернативы по хранилищу}

\begin{table}[H]
\caption{Сравнение подходов к хранению результатов}
\label{tab:storage_alternatives}
\begin{tabularx}{\textwidth}{>{\raggedright\arraybackslash}p{0.24\textwidth}X}
\toprule
Вариант & Комментарий \\
\midrule
NoSQL-only & Гибкость схемы, но менее удобная транзакционная связка для \texttt{events + event\_outbox} \\
PostgreSQL + BLOB-видео & Упрощает единое хранилище, но приводит к раздуванию БД и росту стоимости операций \\
PostgreSQL + файловые артефакты & Транзакционная надежность структурированных данных при контролируемом объеме БД и выносе тяжелых артефактов на файловый слой \\
\bottomrule
\end{tabularx}
\end{table}

В работе выбирается третий вариант: события и конфигурация хранятся в PostgreSQL, а снимки/клипы - в файловом хранилище с сохранением ссылок и метаданных в БД \cite{postgresql_docs}.

\subsection{Архитектурные компромиссы}

Принятый стек формирует набор осознанных компромиссов:
\begin{itemize}
    \item в пользу управляемости: отказ от избыточно сложного распределенного контура на текущем этапе;
    \item в пользу интерпретируемости: правиловой анализ вместо «черного ящика» событийной нейросети;
    \item в пользу надежности: транзакционная таблица исходящих сообщений, повторные попытки и очередь недоставленных сообщений вместо публикации без гарантий доставки;
    \item в пользу воспроизводимости: единый compose-стенд и фиксированные контракты API/сообщений.
\end{itemize}

Такие компромиссы соответствуют логике магистерского исследования: сначала подтвердить корректность целевого контура и измеримые свойства интеграционной эксплуатации на стенде, затем масштабировать отдельные подсистемы.

\section{Исходные данные для разработки и экспериментальной оценки}

Исходные данные для построения и оценки платформы должны отражать реальные условия эксплуатации систем охранного видеонаблюдения и обеспечивать возможность проверки как алгоритмических, так и архитектурных решений.

\subsection{Состав входных данных}

К основным видам исходных данных относятся:
\begin{enumerate}
    \item \textit{Видеопотоки с камер наблюдения}. Данные потоки являются основным источником информации и должны включать сцены с различными условиями освещения, погодными факторами, уровнями плотности объектов, масштабами наблюдения и степенью подвижности сцены.
    \item \textit{Размеченные изображения и видеопоследовательности}. Для обучения и валидации детектора требуются данные с разметкой ограничивающих рамок и классов объектов. Для оценки сопровождения дополнительно необходимы идентификаторы траекторий или эквивалентная разметка перемещения объектов во времени.
    \item \textit{Конфигурация контролируемых зон}. К этой группе относятся полигоны запретных областей, виртуальные линии пересечения, описания периметров, приоритеты сценариев и параметры правил событийного анализа.
    \item \textit{Метаданные камер и потоков}. Необходимы идентификаторы камер, временные метки, характеристики потока, а также сведения о пространственном контексте установки камеры.
    \item \textit{Журналы и эталонные события}. Для оценки модуля событийного анализа требуется набор сценариев или журналов, позволяющих сопоставить сформированные платформой события с ожидаемыми результатами.
\end{enumerate}

Содержательно входные данные должны покрывать не только типовые сценарии, но и сложные случаи, наиболее критичные для задач безопасности: частичные перекрытия объектов, изменение фона, появление мелких объектов на большом расстоянии, резкие изменения освещения, ложные движения, вызванные погодными факторами, а также сцены с несколькими целями одновременно.

\subsection{Требования к подготовке данных}

Для корректной разработки и экспериментальной оценки платформы исходные данные должны удовлетворять ряду требований:
\begin{itemize}
    \item классовый состав разметки должен соответствовать целевым объектам наблюдения;
    \item данные должны включать различные ракурсы, масштабы и режимы съемки;
    \item обучающая и тестовая выборки должны быть разделены таким образом, чтобы оценка качества отражала способность методов к обобщению;
    \item при наличии сценариев с БПЛА требуется отдельное покрытие малыми объектами и сценами дальнего наблюдения;
    \item временная разметка событий должна быть синхронизирована с видеопотоками и траекториями объектов.
\end{itemize}

Наличие структурированных исходных данных особенно важно для раздельной оценки отдельных подсистем платформы: детектора, модуля сопровождения, модуля событийного анализа и механизма передачи сведений о событиях.

\subsection{Показатели оценки качества}

Исходные данные должны обеспечивать возможность вычисления показателей, связанных с целью работы и последующими этапами апробации. К таким показателям относятся:
\begin{itemize}
    \item точность обнаружения объектов, оцениваемая через метрики качества детекции;
    \item устойчивость сопровождения, характеризуемая непрерывностью траекторий, числом переключений идентификаторов и качеством ассоциации;
    \item корректность генерации событий, включая долю пропущенных и ложных событий;
    \item задержка обработки видеопотока и суммарное время от появления объекта до формирования события;
    \item надежность доставки сведений о событиях внешним потребителям.
\end{itemize}

Фиксация указанных показателей на этапе постановки задачи необходима, поскольку именно они позволяют связать выбор методов с целями исследования и впоследствии использовать единый набор критериев при экспериментальной проверке разработанной платформы. При этом показатели p95 \texttt{object->event}, доля отброшенных кадров и метрики событийной точности задаются как целевые показатели расширенной методики; в текущей апробации количественно проверяются интеграционная доставка, отказные сценарии, RBAC-ограничения и демонстрационный БПЛА-контур.

\section{Связь требований с измеряемыми метриками и единой нотацией}

Для обеспечения сквозной сопоставимости главы 1 с главами 2--4 вводится единая нотация входов и выходов конвейера:
\begin{equation}
\mathcal{V} \rightarrow \mathcal{D} \rightarrow \mathcal{T} \rightarrow \mathcal{E} \rightarrow \texttt{event\_outbox}.
\end{equation}
\where
\whereitem{\(\mathcal{V}\)}{видеопотоки}
\whereitem{\(\mathcal{D}\)}{детекции}
\whereitem{\(\mathcal{T}\)}{треки}
\whereitemlast{\(\mathcal{E}\)}{сформированные события}
Такое обозначение далее используется при математической формализации (глава 2), инженерной реализации (глава 3) и в таблицах апробации (глава 4).

Требования главы 1 связываются с метриками следующим образом:
\begin{table}[H]
\caption{Трассировка требований и измеряемых метрик}
\label{tab:req_metrics_traceability}
\begin{tabularx}{\textwidth}{>{\raggedright\arraybackslash}p{0.35\textwidth}X}
\toprule
Требование & Ключевые метрики проверки \\
\midrule
Работа в реальном времени & p95 \texttt{object->event}, длина очередей, доля отброшенных кадров \\
Устойчивость сопровождения & переключения идентификаторов (\(ID\)-switches), доля прерванных траекторий \\
Корректность событий & точность и полнота событий, число ложных/пропущенных событий \\
Надежная доставка & размер очереди исходящих сообщений, число повторных попыток, число сообщений в очереди недоставленных сообщений, доля успешно доставленных событий \\
Безопасность доступа & успешность сценария JWT-аутентификации, матрица RBAC-ограничений \\
\bottomrule
\end{tabularx}
\end{table}

Таким образом, уже на этапе постановки задачи формируется измеримая основа для верификации результатов разработки, что соответствует требованиям воспроизводимости инженерных экспериментов по ГОСТ 7.32 \cite{gost_732_2017} и практикам эксплуатационного мониторинга \cite{prometheus_docs,grafana_docs}.

\section*{Выводы по главе 1}
\addcontentsline{toc}{section}{Выводы по главе 1}

В первой главе выполнен анализ предметной области и сформулирована постановка задачи разработки платформы видеоаналитики реального времени для систем охранного видеонаблюдения. Основные результаты главы могут быть резюмированы следующими положениями.

\textit{Во-первых}, на основе анализа особенностей предметной области показано, что решение рассматриваемой задачи требует согласования четырех уровней обработки: приема видеопотоков, обнаружения объектов в кадре, их сопровождения во времени и событийного анализа траекторий. Сформированные пользовательские роли (\texttt{operator}, \texttt{admin}, \texttt{viewer}) и карты пути пользователя позволили согласовать функциональные требования с реальной деятельностью операторов и администраторов охраняемых объектов.

\textit{Во-вторых}, проведенный обзор существующих платформ видеоаналитики и систематизация литературы по компьютерному зрению, многообъектному сопровождению и микросервисным архитектурам позволили выявить научно-методическую нишу настоящей работы: построение открытого исследовательского решения с фиксированными контрактами событий, прозрачными операционными метриками и воспроизводимым стендом развертывания. Сравнение по таблице~\ref{tab:analogues_comparison} показывает, что существующие промышленные платформы хорошо закрывают типовые сценарии охранной видеоаналитики, однако ограничивают возможности сравнительных исследований и адаптации алгоритмического ядра.

\textit{В-третьих}, на основе функциональных и нефункциональных требований обоснован выбор одноэтапного детектора семейства YOLO, многообъектного модуля сопровождения на основе DeepSORT и правилового подхода к формированию событий, а также выбор технологического стека контура исполнения (FFmpeg/PyAV, FastAPI, PostgreSQL, RabbitMQ, Prometheus, Grafana). Определены состав исходных данных, требования к их подготовке и ключевые показатели качества (p95 \texttt{object} $\rightarrow$ \texttt{event}, переключения идентификаторов, точность и полнота событий, размер очереди исходящих сообщений, число повторов и сообщений в очереди недоставленных), которые далее используются как методическая основа; в главе 4 явно разграничиваются фактически измеренные показатели и показатели, вынесенные в расширенные прогоны.

Совокупность сформированных требований, обоснованного выбора методов и измеримых показателей качества образует техническое задание на разработку платформы и формирует методическую основу для дальнейшего проектирования архитектуры, программной реализации и экспериментальной оценки результатов.
